%% frontmatter/abstract.tex
%% ============================================================================
%% ABSTRACT - A concise summary of your entire thesis
%% ============================================================================
%%
%% STRUCTURE (follow this order):
%% 1. Background/Context (1-2 sentences) - Why is this topic important?
%% 2. Problem Statement (1-2 sentences) - What problem are you solving?
%% 3. Methodology (2-3 sentences) - How did you approach the problem?
%% 4. Key Results (2-3 sentences) - What did you find?
%% 5. Conclusions (1-2 sentences) - What do the results mean?
%% 6. Keywords (5-8 relevant terms)
%%
%% TIPS:
%% - Keep it between 150-300 words
%% - Write in past tense for completed work
%% - Avoid citations, abbreviations, and jargon
%% - Make it self-contained (reader should understand without reading thesis)
%%

\begin{abstract}

% ⚠️ REPLACE THIS SAMPLE TEXT WITH YOUR OWN ABSTRACT

% Background (why is this important?)
[Your research area] is an important field because [reason]. Understanding [specific aspect] is crucial for [application/benefit].

% Problem (what problem are you solving?)
However, current approaches face challenges including [challenge 1], [challenge 2], and [challenge 3]. This research addresses the problem of [specific problem statement].

% Methodology (how did you approach it?)
This study employs [methodology type, e.g., experimental/analytical/simulation-based] approach to investigate [what you investigated]. The research implements [tool/technique/method] using [technology/platform] and evaluates performance through [evaluation method].

% Results (what did you find?)
Experimental results demonstrate that [key finding 1]. Furthermore, [key finding 2]. The proposed approach achieves [quantitative result, e.g., X\% improvement] compared to [baseline/existing methods].

% Conclusion (what does it mean?)
These findings suggest that [implication of results]. The contributions of this research provide [benefit to the field/practitioners].

\vspace{0.5cm}

\textbf{Keywords:} Keyword One, Keyword Two, Keyword Three, Keyword Four, Keyword Five

% TIP: Choose keywords that someone might search for to find your research

\end{abstract}
